% Figure: link-budget waterfall for the geostationary downlink case study.
% Each bar is a running power level in dBW as the signal passes through a stage:
% Tx power, +Tx antenna gain (EIRP), -free-space path loss, +Rx antenna gain,
% giving the received carrier power. The bars are absolute running levels, not
% increments, so the eye reads the level at each stage directly.
\begin{figure}[htbp]
\centering
\begin{tikzpicture}
\begin{axis}[
  width=12cm, height=6.5cm,
  ybar,
  bar width=18pt,
  xlabel={link-budget stage},
  ylabel={running level (dBW)},
  ymin=-160, ymax=60,
  ytick={-160,-120,-80,-40,0,40},
  axis lines=left,
  enlargelimits=0.12,
  symbolic x coords={Tx power, EIRP, after FSPL, at Rx output},
  xtick=data,
  x tick label style={font=\scriptsize, rotate=20, anchor=east},
  nodes near coords,
  nodes near coords style={font=\scriptsize, /pgf/number format/fixed,
    /pgf/number format/precision=0},
]
  % Tx 100 W = 20 dBW; EIRP = 20 + 34 = 54 dBW; after FSPL 205.6 dB = 54 - 205.6
  % = -151.6 dBW; at Rx output add Rx gain 39.7 dB = -111.9 dBW.
  \addplot[engblue, fill=engblue!30] coordinates {
    (Tx power, 20)
    (EIRP, 54)
    (after FSPL, -151.6)
    (at Rx output, -111.9)
  };
\end{axis}
\end{tikzpicture}
\caption[Geostationary downlink budget waterfall]{The running power level
at each stage of the geostationary downlink in the case study, in dBW. The
transmitter delivers $100\,\si{\watt}$ ($20\,\si{\decibel}$ above one watt);
the $34\,\si{\decibel}$ transmit-antenna gain lifts this to an EIRP of
$54\,\si{\decibel}$ above a watt; the $205.6\,\si{\decibel}$ free-space path
loss over the slant range drives the level down to $-151.6\,\si{\decibel}$;
the $39.7\,\si{\decibel}$ receive-dish gain recovers it to a received
carrier power near $-112\,\si{\decibel}$ above a watt. The path loss is the
single dominant term, and the two antenna gains are the engineering levers
that make the link close.}
\label{fig:vol04ch10:link-budget-waterfall}
\end{figure}
